% =====================================================================
%  THÈME 1 — Pondération des équations chimiques
%  Niveau FACILE — Exercices
% =====================================================================
\documentclass[11pt,a4paper]{article}
\usepackage[utf8]{inputenc}
\usepackage[T1]{fontenc}
\usepackage{lmodern}
\usepackage[french]{babel}
\usepackage{amsmath,amssymb,mathtools}
\usepackage{icomma}
\usepackage[version=4]{mhchem}
\usepackage{siunitx}
\sisetup{output-decimal-marker={,},per-mode=symbol}
\usepackage{xcolor}
\usepackage{tikz}
\usepackage[most]{tcolorbox}
\usepackage{enumitem}
\usepackage{array}
\usepackage{fancyhdr}
\usepackage[a4paper,margin=2.1cm,headheight=26pt,headsep=8pt]{geometry}

\definecolor{primary}{RGB}{142,93,168}
\definecolor{primarydark}{RGB}{82,45,108}
\definecolor{excol}{RGB}{0,135,90}

\pagestyle{fancy}\fancyhf{}
\fancyhead[L]{\small\textcolor{primarydark}{\textbf{Fiche 5} — Thème 1 : Pondération}}
\fancyhead[R]{\small\textcolor{excol}{\textsc{Facile}}}
\fancyfoot[C]{\small\thepage}
\renewcommand{\headrulewidth}{0.8pt}
\renewcommand{\headrule}{\hbox to\headwidth{\color{primary}\leaders\hrule height \headrulewidth\hfill}}

\newtcolorbox[auto counter]{exo}[1][]{enhanced,breakable,boxrule=1pt,arc=3pt,
  left=3mm,right=3mm,top=2.4mm,bottom=2.4mm,colback=excol!5,colframe=excol,
  fonttitle=\bfseries,coltitle=white,
  title={Exercice~\thetcbcounter\ \ifx\relax#1\relax\relax\else\textmd{--- \itshape #1}\fi}}

\setlength{\parindent}{0pt}
\setlength{\parskip}{3pt}

\begin{document}

\begin{center}
{\LARGE\bfseries\color{primarydark}Thème 1 — Pondération}\\[1pt]
{\large\color{excol}\textbf{Niveau Facile}}\\[2pt]
{\itshape Objectif : ajuster des coefficients simples et reconnaître le type de réaction.}\\
{\color{primary}\rule{0.6\linewidth}{1pt}}
\end{center}

\begin{exo}[synthèse de l'ammoniac]
Le procédé Haber synthétise l'ammoniac à partir du diazote et du dihydrogène :
\[ \ce{N2 + H2 -> NH3}\qquad\text{(non pondérée).} \]
\textbf{a)} Pondère cette équation avec les plus petits entiers.\\
\textbf{b)} Donne la \textbf{somme} des coefficients stœchiométriques (coefficient~1 compris).
\end{exo}

\begin{exo}[combustion du méthane]
Le méthane brûle complètement dans le dioxygène :
\[ \ce{CH4 + O2 -> CO2 + H2O}\qquad\text{(non pondérée).} \]
\textbf{a)} Pondère l'équation.\quad
\textbf{b)} Combien d'atomes d'oxygène comptes-tu de chaque côté une fois l'équation ajustée ?
\end{exo}

\begin{exo}[oxydation du fer]
Le fer s'oxyde à l'air en oxyde de fer(III) :
\[ \ce{Fe + O2 -> Fe2O3}\qquad\text{(non pondérée).} \]
Pondère l'équation, puis donne la somme des coefficients. \emph{(Indice : commence par le fer,
finis par l'oxygène — cherche le plus petit commun multiple de 2 et 3.)}
\end{exo}

\begin{exo}[repérer une erreur de pondération]
Un étudiant propose, pour la formation de l'alumine :
\[ \ce{2 Al + O2 -> Al2O3}. \]
\textbf{a)} Fais l'inventaire des atomes de chaque côté : l'équation est-elle correctement pondérée ?\\
\textbf{b)} Si non, corrige-la avec les plus petits entiers et donne la somme des coefficients.
\end{exo}

\begin{exo}[nommer le type de réaction]
Sans les pondérer, associe à chacune des réactions suivantes son \textbf{type} parmi :
\textbf{acide--base}, \textbf{redox}, \textbf{précipitation}, \textbf{décomposition}, \textbf{hydrogénation}.
\begin{enumerate}[label=\arabic*., leftmargin=1.8em,topsep=1pt,itemsep=1pt]
  \item \ce{AgNO3 (aq) + KBr (aq) -> AgBr (s) + KNO3 (aq)}
  \item \ce{CaCO3 (s) -> CaO (s) + CO2 (g)}
  \item \ce{C2H4 (g) + H2 (g) -> C2H6 (g)}
  \item \ce{H2SO4 (aq) + 2 NaOH (aq) -> Na2SO4 (aq) + 2 H2O (l)}
  \item \ce{4 Fe (s) + 3 O2 (g) -> 2 Fe2O3 (s)}
\end{enumerate}
\end{exo}

\end{document}
